\sequence{Cross-tabulating two qualitative variables}{Two-way tables}{1--12 March 2027}
\section{Two-way tables and marginal frequencies}
\begin{cours}[Definitions]
A qualitative variable describes a category. It is nominal if its categories have no natural order (means of transport), ordinal if they are ordered (level of satisfaction).
A two-way table counts individuals according to two variables. The sums of the rows and columns are the margins. A marginal frequency divides a margin by the total count.
\end{cours}
\ex{In a population of $100$ individuals, two characteristics $A,B$ are observed:}
\begin{center}\renewcommand{\arraystretch}{1.6}\begin{tabular}{c|rr|r}
&$A$&$\overline A$&Total\\\hline
$B$&20&15&35\\
$\overline B$&12&53&65\\\hline
Total&32&68&100
\end{tabular}\end{center}
Calculate the marginal frequencies of $A$, $B$, $\overline B$, then the frequency of individuals satisfying both $A$ and $B$.
\reponse[2]{$f(A)=32/100=0.32$; $f(B)=35/100=0.35$; $f(\overline B)=65/100=0.65$. The frequency of $A\cap B$ is $20/100=0.20$; it is not a marginal frequency.}
\ex{A Year 12 class of $32$ students is split as follows. Complete the margins and calculate the marginal frequencies ``CS'' and ``girls''.}
\begin{center}\renewcommand{\arraystretch}{1.5}\begin{tabular}{c|rr|r}
&Boys&Girls&Total\\\hline
CS&12&9&\trou{21}\\
No CS&6&5&\trou{11}\\\hline
Total&\trou{18}&\trou{14}&32
\end{tabular}\end{center}
\reponse[1.6]{$f(\mathrm{CS})=21/32=65.625\,\%$; $f(\mathrm{girls})=14/32=43.75\,\%$.}

\fiche{Conditional frequencies and reading sentences}
\begin{cours}[Definition]
The frequency of $A$ \textbf{among the individuals of $B$}, if this subpopulation is nonempty, is
\[f_B(A)=\frac{\text{count satisfying }A\text{ and }B}{\text{count of }B}.\]
The denominator corresponds to the population introduced by ``among''.
\end{cours}
\ex{In the CS table, calculate the frequency of girls among the students without CS, then the frequency of students without CS among the girls. Are the two results equal?}
\reponse[4]{Among the $11$ students without CS, $5$ are girls: $5/11\approx45.45\,\%$. Among the $14$ girls, $5$ do not take CS: $5/14\approx35.71\,\%$. The numerator is the same, the reference populations differ.}
\ex{In the table of $100$ individuals, calculate $f_B(A)$, $f_A(B)$ and $f_{\overline A}(B)$. Translate each result into a sentence.}
\reponse[4]{$f_B(A)=20/35=4/7\approx57.14\,\%$ of the individuals of $B$ satisfy $A$. $f_A(B)=20/32=62.5\,\%$ of the individuals of $A$ satisfy $B$. $f_{\overline A}(B)=15/68\approx22.06\,\%$ of the individuals not satisfying $A$ satisfy $B$.}
\ex{Does ``$60\,\%$ of cyclists are students'' allow us to conclude that $60\,\%$ of students are cyclists?}
\reponse[2.4]{No. For example, out of $10$ cyclists, $6$ are students; if the population has $100$ students in all, only $6\,\%$ of the students are cyclists.}

\fiche{Completing a table from frequencies}
\ex[Fictional situation]{A group has $200$ people, of whom $120$ use a bus ($B$). Among these $120$ people, $75\,\%$ have a pass ($A$). Among the other $80$, $25\,\%$ have a pass.
\begin{enumerate}\item Construct the two-way table $B/\overline B$ and $A/\overline A$ with its margins.\item Calculate the marginal frequency of $A$.\item Calculate the frequency of $B$ among $A$.\end{enumerate}}
\reponse[6]{In $B$, $120\times0.75=90$ people satisfy $A$, so $30$ do not. In $\overline B$, $80\times0.25=20$ satisfy $A$, and $60$ do not.
\begin{center}\begin{tabular}{c|rr|r}&$A$&$\overline A$&Total\\\hline $B$&90&30&120\\$\overline B$&20&60&80\\\hline Total&110&90&200\end{tabular}\end{center}
$f(A)=110/200=55\,\%$. $f_A(B)=90/110=9/11\approx81.82\,\%$.}
\ex{A headline announces ``Three pass holders in four take the bus'', based on $f_B(A)=75\,\%$. Explain the error and propose a correct sentence.}
\reponse[3.2]{The headline reverses the conditioning. The correct sentence is: ``Three people in four among those who take the bus have a pass.'' The proportion of pass holders who take the bus is here about $81.82\,\%$.}

\fiche{Processing a file and collecting real data}
\begin{lstlisting}
# Exemple fictif : (transport, satisfaction)
donnees = [("bus", "oui"), ("marche", "oui"),
           ("bus", "non"), ("velo", "oui"),
           ("bus", "oui"), ("marche", "non")]

bus = [r for r in donnees if r[0] == "bus"]
bus_satisfaits = [r for r in donnees
                 if r[0] == "bus" and r[1] == "oui"]
autres = [r for r in donnees if not r[0] == "bus"]
\end{lstlisting}
\ex{Give the sizes of the three lists. Calculate the frequency of satisfaction among bus users. Replace \texttt{and} by \texttt{or} in the second filter: how many rows are then selected?}
\reponse[3.2]{Sizes: $3,2,3$. Among bus users, the frequency of satisfaction is $2/3$. With \texttt{or}, we select the $3$ bus users and the $2$ satisfied people using another means: $5$ rows, without counting twice the rows satisfying both criteria.}
\begin{methode}[Activity with the class data]
Collect anonymously two variables: main means of transport and journey time in ordered categories ($<15$ min; from $15$ to under $30$ min; $\geq30$ min). One row per response, without names. Construct the two-way table, then two conditional frequencies in opposite directions. State the count, the collection date and any non-responses.
\end{methode}
\ex{Carry out the collection and write two justified conclusions. Explain why the results do not automatically describe all students in Madagascar.}
\reponse[4]{The values depend on the actual collection. Check: one category per variable and per response; sum of the cells equal to the number of usable responses; correct denominator for each frequency; announced rounding. The sample is one particular class, not a representative sample of all students in the country. Non-responses must be distinguished from an observed category.}
